\section{相关工作}

\subsection{元推理与计算行动权衡}

有限理性与元推理研究智能体在计算受限时如何分配内部推理与外部行动\cite{horvitz1987reasoning, russell1991right}。经典单步元推理估计额外一次搜索对动作期望效用的提升，并在增益低于延迟代价时终止搜索\cite{lin2015metareasoning}。在机器人并发控制中，边移动边思考框架权衡继续规划的路径改善与推迟动作切换的跟踪偏差\cite{xiao2020thinking, elboher2023formal}；另一类工作利用历史任务轨迹训练停止策略或预算选择器，由环境特征预测不同算力预算的收益\cite{sung2021quit, magnino2026metareasoning}。近期大模型智能体研究也发现过度思考会挤占交互机会并降低成功率\cite{shen2025thinking, cuadron2025overthinking, hunt2021formalizing}。然而现有方法评估规划收益时普遍直接信任规划器在搜索空间上的自评得分，或依赖大量历史任务训练选择器；本文指出不确定海洋场中自评得分受优化器诅咒污染，并利用留出想象海洋在无历史任务条件下无偏评估计算价值。

\subsection{海洋自适应采样与信息收集}

自主水下航行器在藻华追踪中的应用推动了从规则覆盖向自适应采样的演进\cite{fossum2019}。割草机式全覆盖航线均匀扫测海域\cite{galceran2013coverage}，但在24小时大范围赤潮巡测中仅能覆盖局部水域。信息收集算法常利用高斯过程维护环境均值与方差，采用贪心策略、上置信界采样\cite{srinivas2012gpucb}或汤普森采样\cite{russo2018thompson}驶向高浓度或高不确定性位置，或通过采样树搜索规划多步路径\cite{hollinger2014sampling}。东海原甲藻赤潮遥感研究揭示了其晚春至初夏高发及随潮汐往复漂移的特征\cite{lou2014diurnal, zhao2025distinguishing}。然而现有自适应采样方法均假定机载规划的计算开销可忽略或固定不变，未考虑高功耗芯片与推进电机共享电池及时间的耦合约束；本文将机载算力、芯片功耗与原位复测纳入统一决策空间，刻画了多步规划与机动采样的收益边界。

\subsection{优化器诅咒与集合数据同化}

在随机优化与强化学习中，用同一批有限样本寻优并估计价值会放大随机正误差，即“优化器诅咒”\cite{smith2006curse}；双估计方法通过将动作选择与价值评估解耦至独立样本集来消除过估计\cite{vanhasselt2010double}，而模型误差也会限制有效规划视野\cite{jiang2015horizon}。海洋环境重构方面，业务化水色遥感利用经验正交函数插补生成逐日无云叶绿素场\cite{beckers2003dineof}，结合多尺度融合海表温度\cite{chin2017mur}、水色实测精度评估\cite{ye2021cocts}及联合嵌入预测架构等时空模型\cite{assran2025vjepa2}刻画表层演变，并以序列集合调整卡尔曼滤波\cite{anderson2001eakf}配合五阶分段有理局地化函数\cite{gaspari1999correlation}同化现场观测。目前尚未见将集合同化样本集与双估计思想结合以消除元推理自评偏差的研究；本文将同化更新的想象海洋划分为搜索用海洋与留出海洋，使卫星大数据标定的集合成员充当评判计算价值的独立标尺。
